\chapter{\LaTeX のすゝめ}
\label{ch:examples}

本章は、\LaTeX の基本的な使い方を紹介します。セクション分け、引用、図、表、数式など、必要となる最も一般的な\mynote{そうでもない}要素の実例を示します。

\section{環境}

\subsection{見出し}

\begin{lstlisting}

\part{部}
\chapter{章} 
\section{節}
\subsection{項}
\subsubsection{目}
\paragraph{段落} %段落に見出しをつける場合
\subparagraph{小段落} 

% \chapter*{}、\section*{}、*をつけることで、番号をつけないようにできます。



\end{lstlisting}

\subsection{フォントサイズ}
\LaTeX は、文書内の様々な組版上のニーズに対応するために、フォントサイズコマンドの階層を提供しています。\autorefja{tab:fontsizes} に詳述されている10個の標準フォントサイズコマンドは、脚注や添え字用の \texttt{\textbackslash tiny} から、主要な見出しやタイトル用の \texttt{\textbackslash Huge} まで多岐にわたります。これらのコマンドは、文書の基本フォントサイズ（ドキュメントクラスのオプションで指定）に対する相対的なサイズであり、文書全体を通して一貫したスケーリングを保証します。





\begin{table}[!ht]
\centering

\caption{LaTeX フォントサイズコマンド一覧}
\renewcommand{\arraystretch}{1.8}
% Modern column definition with better spacing
\begin{tabular}{@{\extracolsep{8pt}}l l@{}}
\toprule
\rowcolor{headerblue}
\color{white}\textbf{コマンド} & 
\color{white}\textbf{出力例} \\
\midrule
\rowcolor{lightgray}
\texttt{\textbackslash tiny} & {\tiny This is tiny text} \\
\texttt{\textbackslash scriptsize} & {\scriptsize This is scriptsize text} \\
\rowcolor{lightgray}
\texttt{\textbackslash footnotesize} & {\footnotesize This is footnotesize text} \\
\texttt{\textbackslash small} & {\small This is small text} \\
\rowcolor{lightgray}
\texttt{\textbackslash normalsize} & {\normalsize This is normalsize text} \\
\texttt{\textbackslash large} & {\large This is large text} \\
\rowcolor{lightgray}
\texttt{\textbackslash Large} & {\Large This is Large text} \\
\texttt{\textbackslash LARGE} & {\LARGE This is LARGE text} \\
\rowcolor{lightgray}
\texttt{\textbackslash huge} & {\huge This is huge text} \\
\texttt{\textbackslash Huge} & {\Huge This is Huge text} \\
\bottomrule
\label{tab:fontsizes}
\end{tabular}

\end{table}

Here is a single paragraph that demonstrates the relative scales of standard \LaTeX{} font commands. We start with {\tiny tiny text for fine details}, move up to {\scriptsize script size usually for subscripts}, and then {\footnotesize footnote size}. 

Gradually, we reach {\small small text}, before returning to {\normalsize the default normal size}. To emphasize points, we can use {\large large text}, {\Large larger text for sub-titles}, or {\LARGE even larger text}. Finally, for major impacts, we use {\huge huge} and {\Huge massive sizes}.


これはLaTeXのフォントサイズ変更コマンドを確認するための段落です。まず{\tiny 極小の文字 (tiny)}から始まり、{\scriptsize 添え字などに使うサイズ (scriptsize)}、そして{\footnotesize 脚注サイズ (footnotesize)}へと変化します。

ここから{\small 少し小さな文字 (small)}になり、{\normalsize 標準のサイズ (normalsize)}に戻ります。強調したい場合は{\large やや大きな文字 (large)}、{\Large 大きな文字 (Large)}、あるいは{\LARGE さらに大きな文字 (LARGE)}を使います。見出しなどで目立たせるには{\huge 巨大な文字 (huge)}や、{\Huge 最大級の文字 (Huge)}を使用します。

\subsection{英文の書式設定}

\begin{lstlisting}
% 英文のみ
\textit{}  % イタリック
\textbf{}  % ボールド
\textsf{}  % ゴシック
\end{lstlisting}


The quick brown fox jumps over the lazy dog

\textit{The quick brown fox jumps over the lazy dog}

\textbf{The quick brown fox jumps over the lazy dog}




% Sans-serif regular
\textsf{The quick brown fox jumps over the lazy dog}

% Sans-serif italic
\textsf{\textit{The quick brown fox jumps over the lazy dog}}

% Sans-serif bold
\textsf{\textbf{The quick brown fox jumps over the lazy dog}}



\subsection{略語}
Acronym testing: \ac{UML} -- \acs{UML} -- \acf{UML} -- \acp{UML}

\ac{TPD} -- \acs{TPD} -- \acf{TPD} -- \acp{TPD}

\ac{API}--\acs{DRY}



\subsection{リスト}
箇条書きリストの例:

\begin{lstlisting}
\begin{itemize}
    \item 最初の項目
    \item 2番目の項目
\end{itemize}
\end{lstlisting}

\begin{itemize}
    \item 最初の項目
    \item 2番目の項目
\end{itemize}

\vspace{1cm}
\noindent 番号付きリストの例:

\begin{lstlisting}
\begin{enumerate}
    \item 最初の手順
    \item 2番目の手順
\end{enumerate}
\end{lstlisting}

\begin{enumerate}
    \item 最初の手順
    \item 2番目の手順
\end{enumerate}

\vspace{1cm}
\noindent 説明リストの例:

\begin{lstlisting}
\begin{description}
    \item[項目 1] 最初の手順
    \item[項目 2] 2番目の手順
\end{description}

\end{lstlisting}


\begin{description}
    \item[項目 1] 最初の手順
    \item[項目 2] 2番目の手順
\end{description}


\subsection{Quote}


\begin{lstlisting}
\begin{quote}
    パワー丼食べたい。
\end{quote}
\end{lstlisting}


\begin{quote}
    パワー丼食べたい。    
\end{quote}

\section{参考文献}
本テンプレートでは、APA Style 7the Editionを使用しています。


\begin{center}
    \url{https://apastyle.apa.org/style-grammar-guidelines/references/examples}
\end{center}


日本語は教育心理学研究\mynote{ほぼ}のスタイルです。句読点は英文APAのままです。

\begin{center}
    \url{https://www.jstage.jst.go.jp/browse/jjep/-char/ja/}
\end{center}


\begin{lstlisting}
\parencite{fiorella2022}  % 括弧付き引用
\textcite{fiorella2022}   % 文中引用
\parencite{fiorella2022,knuth1984,vaswani2017}  % 複数の引用
\citeauthor{fiorella2022}  % 著者
\citeyear{fiorella2022}  % 年
\end{lstlisting}


\begin{description}
    
    \item[括弧付き引用:] \parencite{fiorella2022}.
    \item[文中引用:] \textcite{fiorella2022} は...を発見した。
    \item[複数の引用:] \parencite{fiorella2022,knuth1984,vaswani2017}.
    \item[著者と年:] \citeauthor{fiorella2022} の (\citeyear{fiorella2022})
\end{description}





食堂は騒がしかったが、私\mynote{猫}の集中力は一点に注がれていた。目の前にある伝説の料理、東工大パワー丼を見つめていた。多くの者が盲目的にそれを消費する中、私はその構造的完全性を理解しようと試みた。

\textcite{sato2023} が食堂力学に関する独創的な論文で主張したように、水菜と焼肉の正確な配分は単なる料理上の選択ではなく、数学的必然である。私は箸を手に取り、彼らの発見を検証する準備をした。

丼を適切に分析するためには、認知的負荷を最小限に抑える必要があった。私は \citeauthor{mayer2021} の原則を適用した。彼は \citeyear{mayer2021} に、人は言葉だけよりも言葉と絵からより良く学ぶことを確立している。したがって、私は食べる前に丼の写真を撮った。

ご飯を掘り進めると、隠された味のネットワークが疑われた。それはまるで、\textcite{Miede2011} がビジネス関係の文脈で記述した「盗聴」技術のように、どこか背徳的な感じがした。ガーリックソースは豚肉と密かに通信していたのだろうか？

私は隠し味の可能性も考慮した。\textcite{suzuki2024} は最近、バナナに関する包括的なガイドを出版した。パワー丼の中にバナナが隠されている可能性はあるだろうか？迅速な味覚テストの結果、間違いなくそうではないことが確認された。


この食事のエネルギー密度は高い。\textcite{fiorella2022} によって論じられた生成的活動の原則によれば、学習---この場合は消化---は生成的なプロセスである。

しかし、相関関係と因果関係を混同しないように注意しなければならない。\textcite{tanaka2025} がダミー変数に関する研究で警告しているように、豚肉のように見えるものが、実は巧みに偽装されたフライドガーリック（言わば料理上の「ダミー」変数）である可能性がある。


謎は部分的に未解決のままである。\citeauthor{sato2023}の水菜と豚肉の比率に関する理論は真実であったが (\citeyear{sato2023})、パワー丼の感情的インパクトは学術的な引用を超越している。

\section{記号}

正しい入力コードについては \autoref{tab:symbols} を参照してください。コンパイルエラーの最も一般的な原因は、アンパサンド (\&) をエスケープせずに使用することです\mynote{特に.bibファイルの中、著者名のところ}。


\begin{table}[htb]
\centering
\footnotesize
\caption{一般的な LaTeX 記号一覧}
\renewcommand{\arraystretch}{1.5}
% Modern column definition with better spacing
\begin{tabular}{@{\extracolsep{4pt}}l c l c l c@{}}
\toprule
\rowcolor{headerblue}
\color{white}\textbf{入力} & \color{white}\textbf{出力} & 
\color{white}\textbf{入力} & \color{white}\textbf{出力} & 
\color{white}\textbf{入力} & \color{white}\textbf{出力} \\
\midrule
\rowcolor{lightgray}
\texttt{\textbackslash\%} & \% & 
\texttt{\textbackslash\$} & \$ & 
\texttt{\textbackslash\&} & \& \\
\texttt{\textbackslash\{} & \{ & 
\texttt{\textbackslash\}} & \} & 
\texttt{\textbackslash\#} & \# \\
\midrule
\rowcolor{lightgray}
\texttt{\$\textbackslash alpha\$} & $\alpha$ & 
\texttt{\$\textbackslash theta\$} & $\theta$ & 
\texttt{\$\textbackslash pi\$} & $\pi$ \\

\texttt{\$\textbackslash Gamma\$} & $\Gamma$ & 
\texttt{\$\textbackslash Delta\$} & $\Delta$ & 
\texttt{\$\textbackslash Phi\$} & $\Phi$ \\
\bottomrule
\label{tab:symbols}
\end{tabular}
\end{table}

\section{改ページ}

改ページを行うコマンドには、{\textbackslash}pagebreak、{\textbackslash}newpage、{\textbackslash}clearpage があります。それぞれの具体的な使用法は異なりますが、基本的には新しいページを開始するために使用されます。



\section{図}
図は \texttt{gfx/} フォルダに配置してください。

\begin{figure}[htbp]
    \centering
    % Using a standard example image if your specific image isn't available
    \includegraphics[width=0.5\textwidth]{gfx/example_1.jpg} 
    \caption{本館前のイチョウ}
    \label{fig:example-single}
\end{figure}


\begin{figure}[htbp]
    \centering
    \begin{minipage}[b]{0.4\textwidth}
        \centering
        \includegraphics[width=\textwidth]{gfx/example_1.jpg}
        \caption*{(a) 左の図}
    \end{minipage}
    \hfill
    \begin{minipage}[b]{0.4\textwidth}
        \centering
        \includegraphics[width=\textwidth]{gfx/example_2.jpg}
        \caption*{(b) 右の図}
    \end{minipage}
    

    
    \caption{2つの図の全体キャプション}
    \label{fig:two}
\end{figure}

\begin{figure}[htbp]
    \centering
    \begin{minipage}[b]{0.45\textwidth}
        \centering
        \includegraphics[width=\textwidth]{gfx/example_1.jpg}
        \caption*{(a) 1番目の図}
    \end{minipage}
    \hfill
    \begin{minipage}[b]{0.45\textwidth}
        \centering
        \includegraphics[width=\textwidth]{gfx/example_2.jpg}
        \caption*{(b) 2番目の図}
    \end{minipage}
    
    \vspace{1em}
    
    \begin{minipage}[b]{0.45\textwidth}
        \centering
        \includegraphics[width=\textwidth]{gfx/example_3.jpg}
        \caption*{(c) 3番目の図}
    \end{minipage}
    \hfill
    \begin{minipage}[b]{0.45\textwidth}
        \centering
        \includegraphics[width=\textwidth]{gfx/example_4.jpg}
        \caption*{(d) 4番目の図}
    \end{minipage}
    
    \caption{4つの図の全体キャプション}
    \label{fig:all}
\end{figure}


\newpage

\section{表}
表には \texttt{booktabs} を、幅の制御には \texttt{tabularx} を使用してください。

下記のサイトで作れます。

\begin{center}
    \url{https://www.tablesgenerator.com/}
\end{center}



\begin{table}[ht]
    \centering
    \footnotesize
    \caption{booktabs を使用した表の例}
    \label{tab:example}
    \begin{tabularx}{\textwidth}{lX}
        \toprule
        \textbf{列 1} & \textbf{列 2 (可変幅)} \\
        \midrule
        項目 A & 項目 A の説明。長くなる場合は次の行に折り返されます。 \\
        項目 B & 項目 B の説明。 \\
        \bottomrule
    \end{tabularx}
\end{table}


\begin{table}[ht]
\centering
\footnotesize
\caption{booktabs スタイルの電子機器在庫表}

\begin{tabular}{lccr}
\toprule

\textbf{製品} & \textbf{カテゴリ} & \textbf{在庫} & \textbf{価格 (\$)} \\
\midrule
MacBook Pro 16" & Laptop & 45 & 2,499 \\
Dell XPS 15 & Laptop & 32 & 1,799 \\
ThinkPad X1 Carbon & Laptop & 28 & 1,899 \\
\midrule
Magic Mouse 2 & Accessory & 156 & 79 \\
Logitech MX Master 3 & Accessory & 203 & 99 \\
Razer DeathAdder V2 & Accessory & 87 & 69 \\
\midrule
Mechanical Keyboard & Peripheral & 64 & 129 \\
Wireless Keyboard & Peripheral & 91 & 59 \\
Gaming Keyboard RGB & Peripheral & 43 & 159 \\
\midrule
27" 4K Monitor & Display & 38 & 549 \\
34" Ultrawide Monitor & Display & 22 & 899 \\
32" Gaming Monitor & Display & 29 & 699 \\
\midrule
USB-C Hub & Adapter & 245 & 49 \\
Thunderbolt Dock & Adapter & 67 & 279 \\
HDMI Cable 10ft & Cable & 412 & 15 \\
\bottomrule
\end{tabular}
\label{tab:inventory}
\end{table}


\begin{table*}[ht]
  \centering
  \setlength{\tabcolsep}{1.5pt} % default is 6pt

  \caption{長い表の例}
  \label{tab:metrics}
  \footnotesize
  \setlength{\tabcolsep}{2pt}
  \begin{threeparttable}
    % Use S columns defined with Mean(SD) format
    % We need to find the max format for each column across the *whole* table
    % Col 1 (Hist 1 Sim, Sci 1 Sim): Max is 2.2(2.2)
    % Col 2 (Hist 1 Com, Sci 1 Com): Max is 3.2(2.2) (due to '100')
    % Col 3 (Hist 2 Sim, Sci 2 Sim): Max is 2.2(2.2) (due to 15.81)
    % Col 4 (Hist 2 Com, Sci 2 Com): Max is 2.2(2.2)
    \begin{tabular*}{\textwidth}{@{\extracolsep{\fill}}
      l@{\extracolsep{\fill}} 
      S[table-format=2.2(2.2)]
      S[table-format=3.2(2.2)]
      S[table-format=2.2(2.2)]
      S[table-format=2.2(2.2)]
    }
      \toprule
      % History 1 and 2
      & \multicolumn{2}{c}{\textsf{歴史 1}}
      & \multicolumn{2}{c}{\textsf{歴史 2}} \\
      \cmidrule(lr){2-3}\cmidrule(lr){4-5}
      % These are already \multicolumn, so they override 'S' and center
      & \multicolumn{1}{c}{\textsf{単純}} & \multicolumn{1}{c}{\textsf{複雑}}
      & \multicolumn{1}{c}{\textsf{単純}} & \multicolumn{1}{c}{\textsf{複雑}} \\
    \midrule
    %------ Syntactic block ------%
    \addlinespace
    \multicolumn{5}{@{}l}{\textsf{構文的複雑さ}} \\
    文の数 (\emph{N})                  & 27   & 28   & 30   & 29   \\
    文の長さ (トークン)              & 24.44 (12.52) & 46.07 (22.36) & 24.77 (9.06) & 47.83 (23.11) \\
    最大文長                   & 54   & 100  & 43   & 95   \\
    解析木の高さ                     & 4.15 (2.03) & 7.04 (3.01) & 4.43 (1.57) & 7.34 (3.14) \\
    平均依存距離              & 2.62 (0.59) & 3.27 (0.88) & 2.82 (0.40) & 3.19 (0.50) \\

    %------ Lexical block ------%
    \addlinespace
    \multicolumn{5}{@{}l}{\textsf{語彙的複雑さ}} \\
    高頻度トークン (\%) & 17.74 & 14.45 & 17.01 & 13.90 \\
    中頻度トークン (\%)  & 22.91 & 26.38 & 19.50 & 21.07 \\
    低頻度トークン (\%)  &  3.92 &  5.77 &  5.94 &  7.75 \\
    \midrule
    %------ Science 1 and 2 ------%
    & \multicolumn{2}{c}{\textsf{科学 1}}
      & \multicolumn{2}{c}{\textsf{科学 2}} \\
      \cmidrule(lr){2-3}\cmidrule(lr){4-5}
      & \multicolumn{1}{c}{\textsf{単純}} & \multicolumn{1}{c}{\textsf{複雑}}
      & \multicolumn{1}{c}{\textsf{単純}} & \multicolumn{1}{c}{\textsf{複雑}} \\
    \midrule
    %------ Syntactic block ------%
    \addlinespace
    \multicolumn{5}{@{}l}{\textsf{構文的複雑さ}} \\
    文の数 (\emph{N})                  & 26   & 22   & 23   & 25   \\
    文の長さ (トークン)              & 26.31 (7.91) & 57.23 (16.63) & 32.35 (15.81) & 50.32 (19.15) \\
    最大文長                   & 45   & 94   & 87   & 99   \\
    解析木の高さ                     & 5.23 (1.34) & 8.82 (2.30) & 5.61 (1.83) & 8.76 (2.85) \\
    平均依存距離              & 2.90 (0.62) & 2.99 (0.76) & 2.91 (0.51) & 3.12 (0.66) \\

    %------ Lexical block ------%
    \addlinespace
    \multicolumn{5}{@{}l}{\textsf{語彙的複雑さ}} \\
    高頻度トークン (\%) & 14.24 & 13.40 & 14.37 & 16.80 \\
    中頻度トークン (\%)  & 26.36 & 27.05 & 27.75 & 29.30 \\
    低頻度トークン (\%)  &  8.94 &  8.77 &  4.46 &  4.46 \\
    \bottomrule
    \end{tabular*}
    \begin{tablenotes}[flushleft]
      \footnotesize
      \item \textsf{注.} 括弧内の値は標準偏差を表します。ほとんどの指標は、カウントと最大値を除き、文全体の平均を示しています。
    \end{tablenotes}
  \end{threeparttable}
\end{table*}

\clearpage

\section{数式}
数式はインライン $E=mc^2$ または別行立てで記述できます：

\begin{equation}
    \label{eq:example}
    f(x) = \int_{-\infty}^{\infty} \hat{f}(\xi)\,e^{2\pi i \xi x} \,d\xi
\end{equation}

複数行の数式には \texttt{align} を使用します：
\begin{align}
    a &= b + c \\
    &= d + e
\end{align}

\section{相互参照}
章 (\autoref{ch:examples})、図 (\autoref{fig:example-single})、表 (\autoref{tab:example})、数式 (\autoref{eq:example}) を自動的に参照できます。

日本語の場合では、\texttt{autorefja} を使用して、章(\autorefja{ch:examples})、節(\autorefja{sec:bib})、図 (\autorefja{fig:example-single})、表 (\autorefja{tab:example})、数式 (\autorefja{eq:example}) を自動的に参照できます。

あるいは、手動で参照することも可能です：章 (第\ref{ch:examples}章) や節 (\S\ref{sec:otherlanguages})。


\section{bibファイルの書き方}\label{sec:bib}

下記のファイルを参考するしてもよいと思います。
\begin{center}
    \url{https://github.com/ShiroTakeda/jecon-bst/blob/master/jecon-example.bib}
\end{center}

\twemoji{warning}\twemoji{warning}\twemoji{warning}ただし、日本語で「本の一部でそれ自身に表題があるもの」（英文文献の場合incollection）を引用する場合は、\textbf{incollectionja}を使ってください。\twemoji{warning}\twemoji{warning}\twemoji{warning}

下記に例を示す。

\begin{lstlisting}
@InCollectionja{takeda2012_cge,
  author       = {武田, 史郎},
  title        = {応用一般均衡モデルによる地球温暖化対策の分析：有用性と問題点},
  editor       = {有村, 俊秀 and 蓬田, 守弘 and 川瀬, 剛志},
  booktitle    = {地球温暖化対策と国際貿易：排出量取引と国境調整措置をめぐる経済学的・法学的分析},
  chapter      = 1,
  pages        = {15--36},
  publisher    = {東京大学出版会},
  year         = 2012,
  langid       = {japanese},
  yomi         = {takeda, shiro}
}
\end{lstlisting}


\subsection{英語文献}

標準的な BibTeX フォーマットに従って参考文献エントリを作成してください。以下は一般的なエントリタイプの例です：

\begin{lstlisting}

@incollection{fiorella2022,
  author    = {Fiorella, Logan and Mayer, Richard E.},
  title     = {The Generative Activity Principle in Multimedia Learning},
  booktitle = {The Cambridge Handbook of Multimedia Learning},
  editor    = {Mayer, Richard E. and Fiorella, Logan},
  edition   = {3},
  publisher = {Cambridge University Press},
  address   = {Cambridge},
  year      = {2022},
  pages     = {339--350},
  doi       = {10.1017/9781108894333.036}
}

@article{lusato2025,
  author    = {Lu, Jialiang and Sato, Reiko},
  title     = {Linguistic dimensions of comprehensibility and perceived fluency in {L2} speech across tasks of varying complexity},
  journal   = {Journal of Second Language Pronunciation},
  volume    = {11},
  number    = {2},
  year      = {2025},
  pages     = {240--266},
  doi       = {10.1075/jslp.24057.lu}
}

@book{mayer2021,
  author    = {Mayer, Richard E.},
  title     = {Multimedia Learning},
  edition   = {3},
  publisher = {Cambridge University Press},
  address   = {Cambridge},
  year      = {2020},
  doi       = {10.1017/9781316941355}
}

@INPROCEEDINGS{Miede2011,
    author = {Andr{\'e} Miede and G\"{o}khan \c{S}im\c{s}ek and Stefan Schulte
    and Abawi, Daniel F. and Julian Eckert and Ralf Steinmetz},
    title = {{R}evealing {B}usiness {R}elationships -- {E}avesdropping {C}ross-organizational
    {C}ollaboration in the {I}nternet of {S}ervices},
    booktitle = {Proceedings of the Tenth International Conference Wirtschaftsinformatik
    (WI 2011)},
    year = {2011},
    volume = {2},
    pages = {1083--1092},
    isbn = {978-1-4467-9236-0}
}
\end{lstlisting}



\subsection{日本語の文献}


\begin{lstlisting}
@article{tanaka2025,
  title        = {ダミーとミダーの関連性についての検討},
  author       = {田中, 太郎 and 山田, 花子},
  journal      = {日本科学会誌},
  volume       = {10},
  number       = {1},
  pages        = {10--25},
  year         = {2025},
  doi          = {10.1000/dummy.doi.123},
  langid       = {japanese},
  yomi         = {tanaka, taro and yamada, hanako}
}

@book{suzuki2024,
  title        = {バナナの基礎と応用},
  author       = {鈴木, 一郎},
  publisher    = {東工大出版会},
  address      = {東京},
  year         = {2024},
  langid       = {japanese},
  yomi         = {suzuki, ichiro}
}

@inproceedings{sato2023,
  title        = {東工大パワー丼における水菜と豚肉の配分},
  author       = {加藤, 次郎},
  booktitle    = {東工大学食学会第50回全国大会講演論文集},
  pages        = {100--102},
  year         = {2023},
  month        = aug,
  langid       = {japanese},
  yomi         = {kato, jiro}
}
\end{lstlisting}

\texttt{yomi} フィールドの形式によって、参考文献のソート順が決まります：

\begin{description}
    \item[ひらがな:] 日本語文献は英語文献の後に分けてリストされ、五十音順にソートされます。
    \item[ローマ字:] 日本語文献は英語文献と混在し、アルファベット順にソートされます。
\end{description}


中国語の文献でも \texttt{langid=japanese} を使用できますが、参考文献リストでは日本語フォントで表示されます。個人的には問題ないと考えています\twemoji{thumbs up}。




\section{日本語のフォント}

メインフォント\mynote{英文も}はNoto Serif\footnote{a.k.a. 源ノ明朝; Source Han Serif} Japaneseを使っています。

\begin{center}
    \url{https://fonts.google.com/noto/specimen/Noto+Serif+JP}
\end{center}

ゴシックフォント\mynote{英文も}はNoto Sans Japaneseを使っています。

\begin{center}
    \url{https://fonts.google.com/noto/specimen/Noto+Sans+JP}
\end{center}

\subsection{Mainfont}

吾輩は猫である。名前はまだ無い。


\textbf{吾輩は猫である。名前はまだ無い。}

\subsection{Sansfont}


\textsf{吾輩は猫である。名前はまだ無い。}


\textsf{\textbf{吾輩は猫である。名前はまだ無い。}}


\subsection{Monofont}

\texttt{吾輩は猫である。名前はまだ無い。}


\texttt{\textbf{吾輩は猫である。名前はまだ無い。}}


\subsection{特殊記号}

\begin{flushleft}
☀ ☁ ☂ ☃ ★ ☆ ☎ ☏ ☜ ☝ ☞ ☟ ☯ ♀ ♁ ♂ ☉


%% Enclosed Alphanumerics (U+2460-U+24FF)
% Circled numbers 1-20
① ② ③ ④ ⑤ ⑥ ⑦ ⑧ ⑨ ⑩ ⑪ ⑫ ⑬ ⑭ ⑮ ⑯ ⑰ ⑱ ⑲ ⑳

% Circled numbers 21-50 (Enclosed Alphanumeric Supplement)
㉑ ㉒ ㉓ ㉔ ㉕ ㉖ ㉗ ㉘ ㉙ ㉚ ㉛ ㉜ ㉝ ㉞ ㉟

㊱ ㊲ ㊳ ㊴ ㊵ ㊶ ㊷ ㊸ ㊹ ㊺ ㊻ ㊼ ㊽ ㊾ ㊿

% Parenthesized numbers
⑴ ⑵ ⑶ ⑷ ⑸ ⑹ ⑺ ⑻ ⑼ ⑽ ⑾ ⑿ ⒀ ⒁ ⒂ ⒃ ⒄ ⒅ ⒆ ⒇

% Numbers with period
⒈ ⒉ ⒊ ⒋ ⒌ ⒍ ⒎ ⒏ ⒐ ⒑ ⒒ ⒓ ⒔ ⒕ ⒖ ⒗ ⒘ ⒙ ⒚ ⒛

% Parenthesized lowercase letters
⒜ ⒝ ⒞ ⒟ ⒠ ⒡ ⒢ ⒣ ⒤ ⒥ ⒦ ⒧ ⒨ ⒩ ⒪ ⒫ ⒬ ⒭ ⒮ ⒯ ⒰ ⒱ ⒲ ⒳ ⒴ ⒵

% Circled uppercase letters
Ⓐ Ⓑ Ⓒ Ⓓ Ⓔ Ⓕ Ⓖ Ⓗ Ⓘ Ⓙ Ⓚ Ⓛ Ⓜ Ⓝ Ⓞ Ⓟ Ⓠ Ⓡ Ⓢ Ⓣ Ⓤ Ⓥ Ⓦ Ⓧ Ⓨ Ⓩ

% Circled lowercase letters
ⓐ ⓑ ⓒ ⓓ ⓔ ⓕ ⓖ ⓗ ⓘ ⓙ ⓚ ⓛ ⓜ ⓝ ⓞ ⓟ ⓠ ⓡ ⓢ ⓣ ⓤ ⓥ ⓦ ⓧ ⓨ ⓩ



% Double circled numbers
⓵ ⓶ ⓷ ⓸ ⓹ ⓺ ⓻ ⓼ ⓽ ⓾

% Negative circled numbers (black circle)
❶ ❷ ❸ ❹ ❺ ❻ ❼ ❽ ❾ ❿

%% Enclosed CJK Letters and Months (U+3200-U+32FF)

% Parenthesized Ideographs
㈠ ㈡ ㈢ ㈣ ㈤ ㈥ ㈦ ㈧ ㈨ ㈩

㈪ ㈫ ㈬ ㈭ ㈮ ㈯ ㈰

㈱ ㈲ ㈳ ㈴ ㈵ ㈶ ㈷ ㈸ ㈹ ㈺ ㈻ ㈼ ㈽ ㈾ ㈿

㉀ ㉁ ㉂ ㉃


% Circled Katakana
㋐ ㋑ ㋒ ㋓ ㋔ ㋕ ㋖ ㋗ ㋘ ㋙ ㋚ ㋛ ㋜ ㋝ ㋞ ㋟

㋠ ㋡ ㋢ ㋣ ㋤ ㋥ ㋦ ㋧ ㋨ ㋩ ㋪ ㋫ ㋬ ㋭ ㋮ ㋯

㋰ ㋱ ㋲ ㋳ ㋴ ㋵ ㋶ ㋷ ㋸ ㋹ ㋺ ㋻ ㋼ ㋽ ㋾

% Circled Ideographs
㊀ ㊁ ㊂ ㊃ ㊄ ㊅ ㊆ ㊇ ㊈ ㊉

㊊ ㊋ ㊌ ㊍ ㊎ ㊏ ㊐

㊑ ㊒ ㊓ ㊔ ㊕ ㊖ ㊗ ㊘ ㊙ ㊚ ㊛ ㊜ ㊝ ㊞ ㊟

㊠ ㊡ ㊢ ㊣ ㊤ ㊥ ㊦ ㊧ ㊨ ㊩ ㊪ ㊫ ㊬ ㊭ ㊮ ㊯ ㊰

% Japanese Era Names
㋿ ㍻ ㍼ ㍽ ㍾

% Telegraph Symbols
㍿ ㏠ ㏡ ㏢ ㏣ ㏤ ㏥ ㏦ ㏧ ㏨ ㏩ ㏪ ㏫ ㏬ ㏭ ㏮ ㏯

㏰ ㏱ ㏲ ㏳ ㏴ ㏵ ㏶ ㏷ ㏸ ㏹ ㏺ ㏻ ㏼ ㏽ ㏾

%% CJK Compatibility (U+3300-U+33FF) - Square Katakana Words
㌀ ㌁ ㌂ ㌃ ㌄ ㌅ ㌆ ㌇ ㌈ ㌉ ㌊ ㌋ ㌌ ㌍ ㌎ ㌏

㌐ ㌑ ㌒ ㌓ ㌔ ㌕ ㌖ ㌗ ㌘ ㌙ ㌚ ㌛ ㌜ ㌝ ㌞ ㌟

㌠ ㌡ ㌢ ㌣ ㌤ ㌥ ㌦ ㌧ ㌨ ㌩ ㌪ ㌫  ㌭ ㌮ ㌯

㌰ ㌱ ㌲ ㌳ ㌴ ㌵ ㌶ ㌷ ㌸ ㌹ ㌺ ㌻ ㌼ ㌽ ㌾ ㌿

㍀ ㍁ ㍂ ㍃ ㍄ ㍅ ㍆ ㍇ ㍈ ㍉ ㍊ ㍋ ㍌ ㍍ ㍎ ㍏

㍐ ㍑ ㍒ ㍓ ㍔ ㍕ ㍖ ㍗

% Square Time/Date
㍘ ㍙ ㍚ ㍛ ㍜ ㍝ ㍞ ㍟ ㍠ ㍡ ㍢ ㍣ ㍤ ㍥ ㍦ ㍧

㍨ ㍩ ㍪ ㍫ ㍬ ㍭ ㍮ ㍯ ㍰

% Square Latin Abbreviations / Units
㎀ ㎁ ㎂ ㎃ ㎄ ㎅ ㎆ ㎇ ㎈ ㎉ ㎊ ㎋ ㎌ ㎍ ㎎ ㎏

㎐ ㎑ ㎒ ㎓ ㎔ ㎕ ㎖ ㎗ ㎘ ㎙ ㎚ ㎛ ㎜ ㎝ ㎞ ㎟

㎠ ㎡ ㎢ ㎣ ㎤ ㎥ ㎦ ㎧ ㎨ ㎩ ㎪ ㎫ ㎬ ㎭ ㎮ ㎯

㎰ ㎱ ㎲ ㎳ ㎴ ㎵ ㎶ ㎷ ㎸ ㎹ ㎺ ㎻ ㎼ ㎽ ㎾ ㎿

㏀ ㏁ ㏂ ㏃ ㏄ ㏅ ㏆ ㏇ ㏈ ㏉ ㏊ ㏋ ㏌ ㏍ ㏎ ㏏

㏐ ㏑ ㏒ ㏓ ㏔ ㏕ ㏖ ㏗ ㏘ ㏙ ㏚ ㏛ ㏜ ㏝ 
\end{flushleft}



\subsection{Other Languages}\label{sec:otherlanguages}

一応下記の4つの言語のフォントを入れてみました\marginnote{satostyle.sty line 154}。他の言語のフォントを追加すれば、使えるはずです。

例えば、タイ語の場合：

\begin{lstlisting}
% Thai
\newfontfamily\thaifont{Noto Serif Thai}[Script=Thai]
\newcommand{\textthai}[1]{{\thaifont #1}}

\end{lstlisting}


\paragraph{Simplified Chinese}
\textzhsc{海客谈瀛洲，烟涛微茫信难求。 越人语天姥，云霞明灭或可睹。 天姥连天向天横，势拔五岳掩赤城。}

\paragraph{Traditional Chinese}
\textzhtc{海客談瀛洲，煙濤微茫信難求。 越人語天姥，雲霞明滅或可覩。 天姥連天向天橫，勢拔五嶽掩赤城。}

\paragraph{Japanese}
海客瀛洲を談ず、煙濤微茫にして信に求め難しと。越人天姥を語る、雲霓明滅或は睹る可しと。天姥天に連なり天に向って橫たはる、勢は五嶽を拔き赤城を掩ふ。

\paragraph{Korean}
\textkr{
해객담영주 연도미망신난구 월인어천모 운예명멸혹가도 천모연천향천횡 세발오악엄적성}


\section{絵文字}


\newcommand{\smokingMan}{%
\begin{tikzpicture}[baseline=-0.3em, inner sep=0pt]
    \node at (0,0) {\twemoji[height=1cm]{technologist}};
    \node[rotate=15] at (0.8em, 0.05em) {\twemoji[height=0.5cm]{cigarette}};
\end{tikzpicture}
}



論文では使わないと思いますが\smokingMan


\begin{table}[!ht]
    \centering
    \caption{Some Examples}

\begin{tabular}{cll}
\toprule
\textbf{Icon} & \textbf{Command} & \textbf{Description} \\
\midrule
\twemoji{smile} & \verb|\twemoji{smile}| & Standard smile \\
\twemoji{joy} & \verb|\twemoji{joy}| & Tears of joy \\
\twemoji{wink} & \verb|\twemoji{wink}| & Winking face \\
\twemoji{thinking face} & \verb|\twemoji{thinking face}| & Thinking \\
\twemoji{sparkles} & \verb|\twemoji{sparkles}| & Sparkles \\
\bottomrule
\end{tabular}
\end{table}


詳細は下記を参照してください。
\begin{center}
    \url{https://jp.mirrors.cicku.me/ctan/macros/latex/contrib/twemojis/twemojis.pdf}
\end{center}

\section{Margin Note}

\begin{lstlisting}
\mynote{パワー丼食べたい}
\end{lstlisting}

パワー丼食べたい\mynote{その通り}。

\section{ToDo}

下記のように、文中にToDoを追加するのも可能です\todo{学食に行く。}。

\begin{lstlisting}
\todo{パワー丼を食べる。}
\end{lstlisting}

食堂は騒がしかったが、私の集中力は一点に注がれていた。目の前にある伝説の料理、東工大パワー丼を見つめていた\todo{パワー丼を食べる。}。多くの者が盲目的にそれを消費する中、私はその構造的完全性を理解しようと試みた\todo{おかわりする。}。

\section{Draft Mode}

\begin{lstlisting}
\drafttrue   % Draft Modeにある
\draftfalse  % オフにする
\end{lstlisting}

Draft Modeをオンにすると、一番最初にVersion History, 各ページ内に今のVersionと時間が出力されます。\marginnote{In config.tex}

Versionは手動で入力する必要があります。
\begin{lstlisting}
\newcommand{\myVersion}{v0.0.2} 
\end{lstlisting}

\section{余白}

ページのレイアウトは下記の設定で調整できます。\marginnote{In main.tex} \texttt{gutter}と\texttt{fore\_edge}と統一すれば左右の右余白が同じになります。

調整する場合は \texttt{marginnote}と\texttt{todonote}\mynote{使わないなら無視していいです}の余白も合わせて調整する必要があります。

\begin{lstlisting}
\documentclass[
  book,
  twoside,
  openright,
  paper=a4,
  fontsize=11pt, 
  gutter=25mm,      
  fore_edge=40mm,  
  head_space=30mm,
  linegap=0.8zw,  
]{jlreq}
\end{lstlisting}


詳細は下記リンクで確認してください\mynote{わかりずらい}。
\begin{center}
    \url{https://www.tug.org/texlive//Contents/live/texmf-dist/doc/latex/jlreq/jlreq-ja.html}
\end{center}

\section{ルビ}

\ruby{吾輩}{わがはい}は猫である。名前はまだ無い。
どこで生れたかとんと\ruby{見当}{けんとう}がつかぬ。何でも薄暗いじめじめした所でニャーニャー泣いていた事だけは記憶している。吾輩はここで始めて人間というものを見た。しかもあとで聞くとそれは書生という人間中で一番\ruby{獰悪}{どうあく}な種族であったそうだ。この書生というのは時々我々を\ruby{捕}{つかま}えて\ruby{煮}{に}て食うという話である。しかしその当時は何という考もなかったから別段恐しいとも思わなかった。ただ彼の\ruby{掌}{てのひら}に載せられてスーと持ち上げられた時何だかフワフワした感じがあったばかりである。掌の上で少し落ちついて書生の顔を見たのがいわゆる人間というものの\ruby{見始}{みはじめ}であろう。この時妙なものだと思った感じが今でも残っている。第一毛をもって装飾されべきはずの顔がつるつるしてまるで\ruby{薬缶}{やかん}だ。その\ruby{後}{ご}猫にもだいぶ\ruby{逢}{あ}ったがこんな\ruby{片輪}{かたわ}には一度も\ruby{出会}{でく}わした事がない。のみならず顔の真中があまりに突起している。そうしてその穴の中から時々ぷうぷうと\ruby{煙}{けむり}を吹く。どうも\ruby{咽}{む}せぽくて実に弱った。これが人間の飲む\ruby{煙草}{たばこ}というものである事はようやくこの頃知った。

\section{縦書き}

\begingroup

\noindent\hfill % Pushes everything after it to the right
  \vbox{\hsize=5cm \tate 
    % --- Your Vertical Text Here ---
    吾輩は猫である。名前はまだ無い。
    どこで生れたかとんと見当がつかぬ。
    何でも薄暗いじめじめした所でニャーニャー泣いていた事だけは記憶している。
    % -------------------------------
  }

\vspace{1cm}

\noindent\hfill % Pushes everything after it to the right
\vbox{\hsize=5cm \tate 
    % --- Your Vertical Text Here ---
    \ruby{子}{し}\ruby{曰}{いわ}\okurigana{ク}、
    \ruby{学}{まな}\okurigana{ビテ}\kaeriten{レ}
    \ruby{時}{とき}\okurigana{ニ}\ruby{習}{なら}\kaeriten{二}\ruby{之}{これ}\kaeriten{一}\okurigana{ヲ}、
    \ruby{亦}{また}\ruby{説}{よろこ}\okurigana{バ}\kaeriten{レ}
    \ruby{乎}{し}\okurigana{カ}。
}

\endgroup


\noindent\hfill
  \vbox{\hsize=10cm \tate \Large
    
    % ==========================================
    % 1. 訓点文 (Kundoku / Kanbun Style)
    % ==========================================
    \noindent
    {\gtfamily\bfseries 春夜宴桃李園序}
    \vspace{1\zw}


    % 夫(レ)天地(ハ-一)者
    \ruby{夫}{}\okurigana{レ}
    天\ruby{地}{}\okurigana{ハ}\kaeriten{一}者
    % 萬(ばん)物(ぶつ)之逆(げき)旅(りょ)(、)
    \ruby{萬}{ばん}\ruby{物}{ぶつ}
    之\ruby{逆}{げき}\ruby{旅}{りょ}、
    % 光(くわう)陰(いん)(ハ-一)者
    \ruby{光}{くわう}\ruby{陰}{いん}\okurigana{ハ}\kaeriten{一}者
    % 百代之過(くわ)客(かく)(ナリ)(。)
    百代之\ruby{過}{くわ}\ruby{客}{かく}\okurigana{ナリ}。
    % 而(シテ)浮(ふ)生(せい)(ハ)
    \ruby{而}{}\okurigana{シテ}
    \ruby{浮}{ふ}\ruby{生}{せい}\okurigana{ハ} 
    % 若(ごと)(シ-レ)夢(ノ)(。)
    \ruby{若}{ごと}\okurigana{シ}\kaeriten{レ}\ruby{夢}{}\okurigana{ノ}。
    % 爲(スコト-レ)歡(よろこび)(ヲ)
    \ruby{爲}{}\okurigana{スコト}\kaeriten{レ}\ruby{歡}{よろこび}\okurigana{ヲ}
    % 幾(いく)何(ばく)(ゾ)(。)
    \ruby{幾}{いく}\ruby{何}{ばく}\okurigana{ゾ}。
    % 古人秉(とっ)(テ-レ)燭(ヲ)
    古人\ruby{秉}{とっ}\okurigana{テ}\kaeriten{レ}\ruby{燭}{}\okurigana{ヲ}
    % 夜遊(ビシハ)(、)
    夜\ruby{遊}{}\okurigana{ビシハ}、
    % 良(まこと)(ニ)有(ル-レ)以(ゆゑ)(也)
    \ruby{良}{まこと}\okurigana{ニ}\ruby{有}{}\okurigana{ル}\kaeriten{レ}
    \ruby{以}{ゆゑ}\ruby{也}{}。
    % 況(ンヤ)
    \ruby{況}{}\okurigana{ンヤ}
    % 陽春召(ク-二-レ)我(ヲ)
    % Note: Adjusted 'ク二' to Okurigana 'ク' Kaeriten '二'. 
    % Added 'レ' as per your input, though grammatically '二' usually handles the jump.
    陽春\ruby{召}{}\okurigana{ク}\kaeriten{二}
    \ruby{我}{}\okurigana{ヲ}
    % 以(テシ-二)煙景(ヲ-一)(、)
    \ruby{以}{}\okurigana{テシ}\kaeriten{二}煙\ruby{景}{}\okurigana{ヲ}\kaeriten{一}、
    % 大(たい)塊(くわい)
    \ruby{大}{たい}\ruby{塊}{くわい}
    % 假(か)(スニ-レ)我(ニ)
    \ruby{假}{か}\okurigana{スニ}\kaeriten{レ}\ruby{我}{}\okurigana{ニ}
    % 以(テスルヲヤ-二)文章(ヲ-一)(。)
    \ruby{以}{}\okurigana{テスルヲヤ}\kaeriten{二}
    文\ruby{章}{}\okurigana{ヲ}\kaeriten{一}。 
    % 會(くわい)(シテ-二)桃(たう)李(り)之芳(はう)園(ニ-一)(、)
    \ruby{會}{くわい}\okurigana{シテ}\kaeriten{二}
    \ruby{桃}{たう}\ruby{李}{り}
    之\ruby{芳}{はう}\ruby{園}{ゑん}\okurigana{ニ}\kaeriten{一}、
    % 序(じょ)(ス-二)天(てん)倫(りん)之樂(らく)事(ヲ-一)(。)
    \ruby{序}{じょ}\okurigana{ス}\kaeriten{二}
    \ruby{天}{てん}\ruby{倫}{りん}之
    \ruby{樂}{らく}\ruby{事}{じ}\okurigana{ヲ}\kaeriten{一}。
    % 群(ぐん)季(き)(ノ)俊(しゅん)秀(しう)(ナルハ)(、)
    \ruby{群}{ぐん}\ruby{季}{き}\okurigana{ノ}
    \ruby{俊}{しゅん}
    \ruby{秀}{しう}\okurigana{ナルハ}、
    % 皆爲(た)(リ-二)惠(けい)連(れん)(一)(。)
    皆\ruby{爲}{た}\okurigana{リ}\kaeriten{二}
    \ruby{惠}{けい}\ruby{連}{れん}\kaeriten{一}。
    % 吾人(ノ)詠(えい)歌(か)(ハ)(、)
    吾\ruby{人}{}\okurigana{ノ}
    \ruby{詠}{えい}\ruby{歌}{か}\okurigana{ハ}、
    % 獨(り)慚(は)(ズ-二)康(かう)樂(らく)(ニ-一)(。)
    \ruby{獨}{}\okurigana{り}\ruby{慚}{は}\okurigana{ズ}\kaeriten{二}
    \ruby{康}{かう}\ruby{樂}{らく}\okurigana{ニ}\kaeriten{一}。
    % 幽(いう)賞(しやう)未(いま)(ダ-レ)已(や)(マ)(、)
    \ruby{幽}{いう}\ruby{賞}{しやう}
    \ruby{未}{いま}\okurigana{ダ}\kaeriten{レ}\ruby{已}{や}\okurigana{マ}、
    % 高(かう)談(だん)轉(うた)(タ)清(シ)(。)
    \ruby{高}{かう}\ruby{談}{だん}
    \ruby{轉}{うた}\okurigana{タ}\ruby{清}{}\okurigana{シ}。
    % 開(イテ-二)瓊(けい)筵(えん)(ヲ-一)
    \ruby{開}{}\okurigana{イテ}\kaeriten{二}
    \ruby{瓊}{けい}\ruby{筵}{えん}\okurigana{ヲ}\kaeriten{一}
    % 以(テ)坐(シ-レ)花(二)(、)
    \ruby{以}{}\okurigana{テ}
    \ruby{坐}{}\okurigana{シ}\kaeriten{レ}\ruby{花}{}\okurigana{ニ}、
    % 飛(バシテ-二)羽(う)觴(しやう)(ヲ-一)
    \ruby{飛}{}\okurigana{バシテ}\kaeriten{二}
    \ruby{羽}{う}
    \ruby{觴}{しやう}\okurigana{ヲ}\kaeriten{一}
    % 以(テ)醉(フ-レ)月(ニ)(。)
    \ruby{以}{}\okurigana{テ}
    \ruby{醉}{}\okurigana{フ}\kaeriten{レ}\ruby{月}{}\okurigana{ニ}。
    % 不(ンバ-レ)有(ラ-二)佳(か)作(さく)(一)(、)
    \ruby{不}{}\okurigana{ンバ}\kaeriten{レ}
    \ruby{有}{}\okurigana{ラ}\kaeriten{二}
    \ruby{佳}{か}\ruby{作}{さく}\kaeriten{一}、
    % 何(ゾ)伸(ベン-二)雅(が)懐(くわい)(ヲ-一)(。)
    \ruby{何}{}\okurigana{ゾ}\ruby{伸}{}\okurigana{ベン}\kaeriten{二}
    \ruby{雅}{が}\ruby{懐}{くわい}\okurigana{ヲ}\kaeriten{一}。
    % 如(も)(シ)
    \ruby{如}{も}\okurigana{シ}
    % 詩不(ンバ-レ)成(ラ)(、)
    詩\ruby{不}{}\okurigana{ンバ}\kaeriten{レ}\ruby{成}{}\okurigana{ラ}、
    % 罰(ハ)依(ラン-二)金(きん)谷(こく)(ノ)酒(しゅ)數(すう)(ニ-一)(。)
    \ruby{罰}{}\okurigana{ハ}
    \ruby{依}{}\okurigana{ラン}\kaeriten{二}
    \ruby{金}{きん}\ruby{谷}{こく}\okurigana{ノ}
    \ruby{酒}{しゅ}\ruby{數}{すう}\okurigana{ニ}\kaeriten{一}。

    \vspace{1.5\zw}
    
    \small
    李太白「春夜桃李の園に宴するの序」\\
    星川清孝 著「古文真宝選新解」\\
    明治書院 (一九五六) より引用。

  }

\vspace{1cm}

\noindent\hfill
  \vbox{\hsize=10cm \tate \Large
    \noindent
    {\gtfamily\bfseries 春夜桃李の園に宴するの序}
    \vspace{1\zw}



    \ruby{夫}{そ}れ天地は
    \ruby{萬物}{ばんぶつ}の\ruby{逆旅}{げきりょ}、
    \ruby{光陰}{くわういん}は
    百代の\ruby{過客}{くわかく}なり。
    \ruby{而}{しか}して\ruby{浮生}{ふせい}夢の\ruby{若}{ごと}し。
    \ruby{歡}{よろこび}を爲すこと\ruby{幾何}{いくばく}ぞ。
    古人\ruby{燭}{しょく}を
    \ruby{秉}{とっ}て夜遊びしは、
    \ruby{良}{まこと}に\ruby{以}{ゆゑ}有るなり。
    \ruby{況}{いは}んや
    陽春我を\ruby{召}{まね}くに煙景を以てし、
    \ruby{大塊}{たいくわい}我に
    \ruby{假}{か}すに文章を以てするをや。
    \ruby{桃李}{たうり}の\ruby{芳園}{はうゑん}に\ruby{會}{くわい}して、
    \ruby{天倫}{てんりん}の\ruby{樂事}{らくじ}を\ruby{序}{じょ}す。
    \ruby{群季}{ぐんき}の\ruby{俊秀}{しゅんしう}なるは、
    皆\ruby{惠連}{けいれん}たり。
    吾人の\ruby{詠歌}{えいか}は、
    \ruby{獨}{ひと}り\ruby{康楽}{かうらく}に\ruby{慚}{は}ず。
    \ruby{幽賞}{いうしやう}
    \ruby{未}{いま}だ\ruby{已}{や}まず。
    \ruby{高談}{かうだん}\ruby{転}{うた}た清し。
    \ruby{瓊筵}{けいえん}を開いて以て花に坐し、
    \ruby{羽觴}{うしやう}を飛ばして以て月に醉ふ。
    佳作有らずんば、何ぞ\ruby{雅懐}{がくわい}を伸べん。
    \ruby{如}{も}し詩成らずんば、罰は\ruby{金谷}{きんこく}の
    \ruby{酒數}{しゅすう}に\ruby{依}{よ}らん。


}
  
\section{main.texの構造と使い方}

この \texttt{main.tex} ファイルは、論文全体の構成をコントロールするものです。

実際の文章（章ごとの内容）はすべて別のファイル（\texttt{Chapters/} フォルダなど）に分かれており、この \texttt{main.tex} はそれらを\textbf{「どの順番で、どのように並べるか」}を指示する役割を持っています。

以下に、各ブロックの役割と使い方を解説します。

\subsection{前付（Front Matter）}
論文の本文に入る前の、表紙や目次などです。

\begin{lstlisting}
% --- TITLE PAGE ---
\include{FrontBackmatter/Titlepage}       % 表紙
\include{FrontBackmatter/Titleback}       % 表紙の裏
\cleardoublepage\include{FrontBackmatter/Committee}     % 審査員一覧
\cleardoublepage\include{FrontBackmatter/Dedication}    % Dedication
\cleardoublepage\include{FrontBackmatter/Abstract}      % アブストラクト
\cleardoublepage%*******************************************************
% Publications
%*******************************************************
\pdfbookmark[1]{業績}{publications}
\chapter*{業績}

\section*{学術論文}
\subsection*{査読付き}
\noindent\hangindent=2em
Author, A. (Year). Title of article. \textit{Journal Name, Volume}(Issue), Pages. \url{https://doi.org/...}

\noindent\hspace{2em}\textit{This research forms the foundation of Chapter X.}

\vspace{0.1in}

\section*{学会発表}

\subsection*{査読付き}

\noindent\hangindent=2em
Author, A. (Year, Month Date). \textit{Title of presentation} [Poster/Oral presentation]. Conference Name, Location.

\noindent\hspace{2em}\textit{This research forms the foundation of Chapter Y.}

\subsection*{査読なし}

\noindent\hangindent=2em
Author, A. (Year, Month Date). \textit{Title of presentation} [Poster/Oral presentation]. Conference Name, Location.

\noindent\hspace{2em}\textit{This research forms the foundation of Chapter Y.}  % 研究業績リスト
\cleardoublepage%*******************************************************
% Acknowledgments
%*******************************************************
\pdfbookmark[1]{謝辞}{acknowledgments}



\bigskip

\begingroup
\let\clearpage\relax
\let\cleardoublepage\relax
\let\cleardoublepage\relax
\chapter*{謝辞}

Put your acknowledgments here.

...


\endgroup
 % 謝辞
\cleardoublepage%*******************************************************
% Table of Contents
%*******************************************************
% NOTE: Page style is already set to 'plain' in main.tex frontmatter section
% DO NOT override it here or it will persist into mainmatter!
\refstepcounter{dummy}
\pdfbookmark[1]{\contentsname}{tableofcontents}
\setcounter{tocdepth}{2} % toc depth
\setcounter{secnumdepth}{3} % section numbering depth
\tableofcontents
\markboth{\contentsname}{\contentsname}
%*******************************************************
% List of Figures and Tables
%*******************************************************
%\cleardoublepage
\newpage
\refstepcounter{dummy}
\pdfbookmark[1]{\listfigurename}{lof}
\listoffigures

\vspace{2em} % Space between lists if on same page
\refstepcounter{dummy}
\pdfbookmark[1]{\listtablename}{lot}
\listoftables


\pdfbookmark[1]{略語一覧}{acronyms}
\markboth{Acronyms}{Acronyms}
\chapter*{略語一覧}
\renewcommand{\aclabelfont}[1]{\small\HeadSansRegular{#1}}
\setlength{\labelsep}{1em} % Increase space between label and description
\begin{acronym}[WWWW]
    \acro{DRY}{Don't Repeat Yourself}
    \acro{API}{Application Programming Interface}
    \acro{UML}{Unified Modeling Language}
    \acro{TPD}{Tokyo-Tech Power Don}
\end{acronym}      % 目次
\end{lstlisting}

\begin{itemize}
    \item \textbf{使い方}: \texttt{FrontBackmatter} フォルダの中にある各 \texttt{.tex} ファイル\marginnote{Contents.tex触らないほうがいい}を編集してください。
    \item \textbf{\texttt{\textbackslash cleardoublepage}}: 「奇数ページ（右側のページ）から始める」という命令です。両面印刷の際、新しいセクションが必ず右ページから始まるように調整します。
\end{itemize}

\subsection{本文（Main Matter）}
ここからページ番号が「アラビア数字（1, 2, 3...）」に切り替わります。

\begin{lstlisting}
% --- MAINMATTER: Fancy style ---
\mainmatter          % ページ番号を「1」にリセット
\pagestyle{fancy}    % ヘッダーを表示する設定
\end{lstlisting}

\subsubsection{部（Part）と章（Chapter）の構成}
このテンプレートでは、複数の「章」をまとめて「部（Part）」として整理できるスタイルになっています。

\begin{lstlisting}
\cleardoublepage
% 部（Part）の扉ページに載せる説明文
\ctparttext{This part provides essential information...}

% 「第I部 ガイド」という扉ページを作成
\part{ガイド}\label{pt:guide}

% 実際の章ファイルを読み込み
\chapter{はじめに}\label{ch:intro}



論文の組版において \LaTeX やWord\mynote{Wordは苦手です。}と格闘しておられる方も多いのではないでしょうか。
このテンプレートは、
\texttt{classicthesis} パッケージ  4.8 (\url{https://ctan.org/pkg/classicthesis?lang=en}) をベースに、我々の執筆ニーズに合わせてカスタマイズしたものです。

主な変更点は以下の通りです：
\begin{itemize}
    \item APA 第7版の引用スタイル
    \item 日本語著者名のカスタムフォーマット
    \item 日本語文献のソート機能
    \item 多言語サポート
\end{itemize}


\autorefja{ch:examples}と\autorefja{ch:cat}では、実例を用いて様々な \LaTeX の機能を紹介します。

\section{クイックスタート}

\textbf{すぐに執筆を始めたいですか？} 基本的な手順は以下の通りです：

\begin{enumerate}
    \item \textsf{情報の入力}: \texttt{config.tex} (14行目以降) を編集し、名前、タイトル、所属などを入力します。

    \item \textsf{執筆}: \texttt{Chapters/Chapter01.tex}、\texttt{Chapter02.tex} などに章ごとのコンテンツを記述します。

    \item \textsf{文献の構築}: \texttt{Bibliography.bib} (または \texttt{part1.bib}、\texttt{part2.bib} など) に参考文献エントリを追加します。

    \item \textsf{画像の追加}: 図表を \texttt{gfx/} フォルダに配置します。

    \item \textsf{コンパイル}: メニュー (左上) を開き、コンパイラを LuaLaTeX に設定します。

    \item \textsf{Configuration Options}: In  texttt{config.tex}, you can enable specific settings. Enable \texttt{drafting} to print the version information on the first page of the thesis. Or set the metadata.
    
    \item \textsf{Version History}: In \texttt{VersionHistory.tex}, change the version history details.
\end{enumerate}



\section{英文論文の場合}

英文がメイン場合は下記のテンプレートを使用してください。

\begin{center}
\url{https://github.com/0ldriku/ClassicThesis_EN}
    
\end{center}

\cleardoublepage
\end{lstlisting}

\begin{description}
    \item[\texttt{\textbackslash ctparttext\{...\}}] 部（Part）のタイトルページに、その部の概要説明を入れることができます。
    \item[\texttt{\textbackslash part\{...\}}] 論文を大きく区切るための「部」を作ります。
    \item[\texttt{\textbackslash include\{...\}}] 実際に執筆する章ファイル（\texttt{.tex}）をここで読み込みます。新しい章を増やす場合は、ここに追記します。
\end{description}

\subsection{付録（Appendix）と参考文献}
本文が終わった後の資料セクションです。

\begin{lstlisting}
\appendix             % 番号体系が「付録A」などに変わります
\cleardoublepage
\part{付録}
\include{Chapters/AppendixA}
\include{Chapters/AppendixB}

% 目次の見た目を整える
\addtocontents{toc}{\protect\addvspace{0.5pc}}
\cleardoublepage

% 参考文献リストを出力
\printbibliography[heading=bibintoc, title={引用文献}]
\end{lstlisting}

\begin{itemize}
    \item 付録が必要なら \texttt{Chapters/AppendixA.tex} などを編集します。
    \item 参考文献は \texttt{.bib} ファイルに基づき自動生成されます。
\end{itemize}

\subsection{執筆のアドバイス}

\begin{enumerate}
    \item \textbf{\texttt{main.tex} は編集しない}: 主に章の順番を入れ替えたり、新しい章を追加するときだけ編集します。
    \item \textbf{執筆は個別のファイルで}: 文章を書くときは、\texttt{Chapters/Chapter01.tex} などの個別のファイルを開いて書いてください。
    \item \textbf{コンパイルは常に \texttt{main.tex} で}: どのファイルを編集していても、コンパイルの対象はこの \texttt{main.tex} になるように設定してください（OverleafのMain document設定やマジックコメントを活用）。
\end{enumerate}


